\chapter{[LER] Lernen lernen}

\section{Literatur}

\section{Prüfungen}

\subsection{Anweisungen}

\begin{enumerate}
  \item Am Begin eines Fallbeispieles erhalten Sie die Informationen zur
  Situation und Szene.
  \item Dokumentieren Sei jeden Befund (laut ansagen)! Erst dann passt der
  Prüfer den Befund an das Prüfungsbeispiel an.
  \item Die simulierte Funkkennung lautet \TakeNotesMedium . 
  \item Telefongespräche, z.B. mit der Leitstelle, sind mit dem simulierten
  Telefon zu führen. Die gewählte Rufnummer ist laut anzusagen.
\end{enumerate}

\subsection{Globalskalen}

\begin{gWideParEnv}

\begin{tabulary}{\linewidth}{LLLLL}
1 & 2 & 3 & 4 & 5
\\\addlinespace\hline\addlinespace
\multicolumn{5}{c}{\TabHeading{Untersuchungsablauf}}
\\\addlinespace
maximale Effizienz im Arbeitsablauf,klare Bewegungsökonomie (müheloser Fluss
von einer Bewegung zur nächsten) 

Durchführung weitgehend vollständig\,/\,in weitgehend korrekter Reihenfolge
&
effizienter Arbeitsablauf überwiegt, vereinzelt unnötige Bewegungen

Durchführung weitgehend vollständig\,/\,in weitgehend korrekter Reihenfolge
durchgeführt, auf Widrigkeiten im Ablauf wird korrekt reagiert
&
teilweise effizienter Arbeitsablauf, vermehrt unnötige Bewegungen

Durchführung ausreichend vollständig\,/\,in ausreichend korrekter Reihenfolge
durchgeführt; auf Widrigkeiten im Ablauf wird nur teilweise korrekt reagiert
&
Ineffizienter Arbeitsablauf, viele unnötige Bewegungen und/oder Unsicherheiten
bezüglich der nächsten Bewegung

Durchführung in relevanten Teilbereichen unvollständig oder in falscher
Reihenfolge; auf Widrigkeiten im Ablauf kann nicht oder nur unangemessen
reagiert werden
&
Durchführung patienten-, eigen- oder drittgefährdend oder für das
Patienten-Outcome wesentliche Untersuchungen wurden nicht durchgeführt
\\\addlinespace\hline\addlinespace
\multicolumn{5}{c}{\TabHeading{Instrumentenhandhabung}}
\\\addlinespace
Flüssige, routinierte Bewegungen mit den Instrumenten

gschickte Instrumentenanwendung
&
kompetente Verwendung der Instrumente überwiegt

vereinzelt ungeschickte Instrumentenanwendung
&
unroutinierte Verwendung der Instrumente

vermehrt ungeschickte Instrumtenenanwendung
&
unangemessene Verwendung der Instrumente

deutlich ungeschickte Instrumentenanwendung
&
patienten-, eigen- oder drittgefährdend
\\\addlinespace\hline\addlinespace
\multicolumn{5}{c}{\TabHeading{Behandlung und Maßnahmen}}
\\
Optimale Behandlung

Durchführung weitgehend vollständig\,/\,in weitgehend korrekter Reihenfolge
&
kompetente Behandlung überwiegt, vereinzelt unnötige Bewegungen

Behandlung weitgehend vollständig\,/\,in weitgehend korrekter Reihenfolge
durchgeführt, auf Widrigkeiten wird korrekt reagiert
&
teilweise kompetente Behandlung, teilweise unnötige Maßnahmen

Behandlung ausreichend vollständig\,/\,in ausreichend korrekter Reihenfolge
durchgeführt; auf Widrigkeiten wird nur teilweise korrekt reagiert
&
Ineffiziente Behandlung oder falsche Behandlung \E{ohne Einfluss auf das
Patienten-Outcome}, viele unnötige Maßnahmen und/oder Unsicherheiten bezüglich
der nächsten Maßnahme

Durchführung in relevanten Teilbereichen unvollständig oder in falscher
Reihenfolge; auf Widrigkeiten kann nicht oder nur unangemessen
reagiert werden
&
Durchführung patienten-, eigen- oder drittgefährdend oder für das Patienten-Outcome
wesentliche  Maßnahmen
wurden nicht durchgeführt 
\\\addlinespace\hline\addlinespace
\multicolumn{5}{c}{\TabHeading{Haltung gegenüber dem Patienten}}
\\\addlinespace
aufmerksam, warm, entspannt, umgänglich, geht auf die Bedürfnisse des Pateinten
ein

optimales Vorstellen, optimale Erklärung der Prozedur, optimales Verabschieden
&
Angemessene, aber eher neutrale Grundhaltung

Korrektes Vorstellen, Korrekte Erklärung der Prozedur, Korrektes Verabschieden
&
Unbeholfen oder unsicher oder zögerlich oder überheblich oder kühl oder
gehetzt oder abrupt

Mängel beim Vorstellen und/oder bei der Erklärung der Prozedur und/oder beim
Verabschieden
&
unpersönlich, geht nicht auf die Bedürfnisse des Patienten ein

kein oder unzureichendes Vorstellen, keine oder unzureichende Erklärung der
Prozedur, kein oder unzureuchendes Verabschieden
&
beleidigend oder prvozierendes Auftreten oder patienten-, eigen- oder
drittgefährdend 
\\\addlinespace\hline\addlinespace
\multicolumn{5}{c}{\TabHeading{Gesprächsverlauf}}
\\
maximale Effizienz im Einholen der Informationen

Idealer Wechsel zwischen offenem und geschlossenem Frgatentyp

strukturierter, aber flexibler Gesprächsfluss
&
Überwiegend Effizienz iom Einholen der Informationen

angemessener Wechsel zwischen offenem und geschlossenem Fragentyp

überwiegend strukturierter, aber flexibler Gesprächsfluss
&
teilweise Effizenz im Einholen der Informationen

kein Wechsel des Fragentyps und/oder unangemesenes Verwenden
eines Fragentyps

eher unstrukturierter und/oder starrer Gesrächsfluss
&
Ineffizienz im Einholen der Information (z.B. Unsicherheit über die als
nächstes zu stellende Frage)

Unangemessener Fragentyp überwiegt (Suggestivfragen)

unstrukturierter und/oder starrer Gesprächsfluss
&
erheben von falschen Informationen z.B. durch Suggestivfragen

patienten-, eigen- oder drittgefährdend oder für das Patienten-Outcome
wesentliche Informationen wurden nicht eingeholt
 \\\addlinespace\hline
\end{tabulary}


\end{gWideParEnv}

